\section{Related Work}

\subsection{Research Background}

Path planning is a cornerstone technology in the evolution of robotic systems. In warehouse environments, path planning algorithms must supply automated guided vehicles (AGVs) with collision-free trajectories while satisfying optimization criteria such as energy consumption, transport time, and communication latency \cite{ref1}. The objective of AGV path planning research is to navigate an AGV from its origin to its destination in as few steps as possible, avoiding all obstacles along the way, thereby improving the overall performance of the warehouse system.

\subsubsection{Traditional Path Planning Algorithms}

A substantial body of research has addressed this domain, yielding a variety of traditional algorithms that broadly fall into two categories: global path planning and local path planning \cite{ref2}.

Global path planning algorithms assume that the robot possesses complete knowledge of the obstacle layout in its surroundings. These methods first construct an abstract map model of the global environment, then perform a global search over this model \cite{ref3} to derive an optimal or near-optimal path that safely guides the robot to its target. Representative global path planning algorithms include graph search methods, sampling-based methods, and bio-inspired intelligent algorithms. Graph search methods, such as Dijkstra's algorithm \cite{ref4} and A* \cite{ref5}, compute a feasible path from start to goal based on known map information; however, the state-space search paradigm underlying these methods incurs poor computational efficiency in complex scenes due to the large volume of data and computation time required. Sampling-based planning algorithms, with the Rapidly-exploring Random Tree (RRT) algorithm as a prominent example, are particularly well-suited to single-AGV path planning in static environments, yet they also contend with low search efficiency and susceptibility to local optima. Bio-inspired intelligent algorithms, including genetic algorithms, particle swarm optimization, and ant colony optimization, solve path planning problems by simulating natural phenomena such as biological evolution, foraging behavior, and nest construction. These methods search for optimal solutions by mimicking individual evolution and collective swarm behavior, which makes them suitable for multi-robot systems and yields relatively high solution efficiency. In practice, however, they are prone to converging to local optima and exhibit slow convergence speed.

Local path planning algorithms are designed for scenarios in which the robot can perceive only partial environmental information through its sensors. The robot must autonomously plan a collision-free optimal path in real time based on the locally sensed data. Owing to their reliance on local environmental features, these methods may produce only locally optimal paths rather than globally optimal ones, and in some cases may fail to reach the destination altogether \cite{ref6}. Widely adopted local path planning algorithms include the artificial potential field (APF) method \cite{ref7} and the dynamic window approach (DWA) \cite{ref8}. The APF method models the robot as moving within a virtual force field, where obstacles exert repulsive forces and the target exerts an attractive force; the resultant of these forces guides the robot through the environment. While computationally simple and capable of real-time obstacle avoidance, the APF method is susceptible to local minima and may exhibit oscillatory behavior in narrow passages \cite{ref9}. The DWA is a velocity-based local planning method that imposes velocity constraints to satisfy kinematic and obstacle-avoidance requirements, searching for the optimal control velocity in the velocity space to achieve swift and safe navigation. Its success rate, however, depends heavily on global parameters, and it is prone to failure in unfamiliar environments.

Traditional algorithms have been widely applied to warehouse robot path planning. Wang et al.\ \cite{ref10} proposed a complex diagonal distance metric to address the shortcomings of conventional distance functions in A*, reducing both the number of searched nodes and the planned path length. Li et al.\ \cite{ref11} introduced PQ-RRT*, a novel algorithm that achieves both asymptotic optimality and rapid convergence, overcoming the slow convergence of the original RRT. Liu et al.\ \cite{ref12} developed an improved ant colony algorithm incorporating an adaptive pheromone evaporation coefficient strategy and a load-balancing strategy, which enhanced the algorithm's search efficiency. Wang et al.\ \cite{ref13} constructed virtual target points on top of the original APF method and combined them with a fuzzy control approach to design a finite state machine, improving the robot's adaptability in complex environments. Xiao et al.\ \cite{ref14} proposed the IACO-DWA algorithm for robot path planning in large-scale complex environments, refining the evaluation function of the traditional DWA to further improve path safety and smoothness.

In summary, while traditional path planning algorithms have achieved considerable research progress, they remain subject to fundamental limitations. They typically rely on pre-constructed environmental map models, which compromises their effectiveness in dynamic and complex settings. Furthermore, most of these methods are designed to plan paths from a fixed start point to a fixed goal point, rendering them incapable of accommodating real-time changes in the target destination. This means the robot is effectively restricted to a locally optimal solution, making these algorithms ill-suited to the dynamically changing environments and tasks characteristic of intelligent warehouses.

\subsubsection{Reinforcement Learning-Based Path Planning Algorithms}

Reinforcement learning (RL) is an emerging direction in artificial intelligence research. Through interaction with the environment, an RL agent receives feedback signals in the form of rewards or returns and makes a sequence of decisions aimed at maximizing cumulative reward. Owing to their independence from pre-specified environmental map models and prior knowledge, RL algorithms have increasingly become a research focus for path planning in various complex environments.

Notable progress has been made in applying RL to path planning. Soong et al.\ \cite{ref21} proposed a partially guided Q-learning algorithm that employs the flower pollination algorithm (FPA) to improve the initialization process, addressing the problem of slow convergence to the optimal solution. Zhao et al.\ \cite{ref22} introduced an experience-memory Q-learning algorithm that continuously updates the shortest distance from the current state node to the starting point; by using an EM table to record distance information and a Q-table to assist in guiding experience transfer and experience reuse strategy updates, the algorithm improves the robot's learning efficiency. Xu et al.\ \cite{ref23} proposed an improved Q-learning algorithm that incorporates a gravitational potential field during Q-value initialization and adds directional factors to the state set; experimental results demonstrate that the improved algorithm requires less planning time and produces smoother paths. Wang \cite{ref24} integrated the heuristic function of the A* algorithm into the action selection strategy of conventional Q-learning, thereby accelerating convergence speed.

To address the pervasive issue of sparse rewards in RL, Andrychowicz et al.\ introduced the Hindsight Experience Replay (HER) mechanism. The core idea of HER is reward relabeling: states that the agent actually reaches, even if they are not the intended goal, are retroactively designated as virtual goals, thereby extracting useful learning signals from failed experiences. Although HER performs well with static goals, extending its relabeling paradigm to topologically dynamic environments where obstacles change over time remains a challenging open problem.

\subsection{Summary and Problem Analysis}

A synthesis of the above research on traditional path planning algorithms and reinforcement learning in the domain of AGV scheduling reveals that, while existing methods have achieved meaningful results in specific scenarios, they still face intractable theoretical and engineering bottlenecks when confronted with the highly dynamic, high-density characteristics of Robotic Mobile Fulfillment Systems (RMFS). The limitations of current research concentrate on the following three fronts.

First, there is a fundamental conflict between the optimality guarantees of traditional heuristic algorithms and the overhead of dynamic replanning. Traditional global path planning algorithms such as A* and Dijkstra can reliably compute collision-free globally optimal paths in known static environments. RMFS scenarios, however, are characterized by extreme dynamic complexity: AGVs inevitably encounter randomly moving pedestrians, other operating AGVs, or temporarily scattered materials during execution. Under such highly dynamic conditions, once the environmental topology changes due to emergent obstacles, traditional graph search algorithms must perform high-frequency replanning for the affected nodes or even the entire global map. As the system scale expands and AGV density increases, the computational overhead of such replanning grows exponentially, engendering severe communication and decision-making delays that substantially constrain the overall throughput of the warehouse system.

Second, standard RL suffers from acute sparse reward and convergence difficulties under complex constraints. Deep reinforcement learning, with its paradigm of offline training and online decision-making, offers a promising alternative for dynamic obstacle avoidance. Algorithms such as standard DQN can produce action commands at the millisecond level, yet they exhibit fatal shortcomings when deployed in practical RMFS settings. RMFS imposes stringent physical operating rules, for instance that fully loaded AGVs are prohibited from traversing storage stations and non-target picking stations. These rules cause the state-action space to be saturated with illegal actions. During the prolonged early exploration phase, it is exceedingly difficult for the agent to trigger the terminal reward of task completion through random exploration alone. This severe sparse reward problem, compounded by an extremely high illegal action rate, leads to the accumulation of massive amounts of ineffective noisy data in the experience replay buffer. The neural network is unable to extract effective policy gradients from such data, ultimately resulting in extremely slow convergence or outright failure of the model.

Third, existing reward shaping mechanisms exhibit clear limitations in dynamic topological scenarios. To mitigate the sparse reward problem in RL, the research community widely adopts reward shaping mechanisms to provide dense learning signals. Mainstream studies, however, predominantly rely on static potential functions, such as those derived from the artificial potential field method or the heuristic distance of a static A* planner. In RMFS environments, the random movement of dynamic obstacles can sever the original optimal topological corridor at any moment. If static potential-based reward shaping is retained under these conditions, the static potential field not only fails to provide correct obstacle-avoidance guidance when the optimal path is temporarily blocked, but also persistently outputs erroneous gradient signals that direct the agent to traverse obstacles. Such misleading rewards can cause the AGV to oscillate locally near obstacle boundaries, enter deadlock, or even increase the risk of collision.

Taken together, these observations indicate that the field of AGV path planning for complex dynamic RMFS scenarios still lacks an effective integrative framework. Such a framework must simultaneously preserve the global optimal guidance capability of classical heuristic algorithms in complex topological structures and fully harness the millisecond-level online responsiveness of deep reinforcement learning when confronting emergent dynamic obstacles.

To address this gap, this paper proposes a single-AGV path planning algorithm termed DRS-DQN (D* Lite Reward Shaping-DQN), which is guided by the dynamic potential of the D* Lite distance field. Chapter~3 elaborates on how this algorithm achieves obstacle avoidance and stable convergence in dynamic warehouse environments through backtracking reward shaping and a deferred computation mechanism, while ensuring policy optimality.
